\documentclass[11pt,a4paper]{article}

\usepackage[T1]{fontenc}
\usepackage[utf8]{inputenc}
\usepackage[slovak]{babel}
\usepackage{lmodern}
\usepackage[margin=2.2cm]{geometry}
\usepackage{microtype}
\usepackage{array}
\usepackage{tabularx}
\usepackage{booktabs}
\usepackage{enumitem}
\usepackage{xcolor}
\usepackage{amsmath}
\usepackage{tikz}
\usepackage{fancyhdr}
\usepackage{lastpage}
\usepackage{cite}
\usepackage{url}
\usepackage[hidelinks]{hyperref}

\hypersetup{
  pdftitle={Black Box -- VV-F vecný zámer projektu},
  pdfauthor={Kyryl Tkachenko, Tymofii Trembovetskyi, Alex Vyšný}
}

\setlength{\parindent}{0pt}
\setlength{\parskip}{0.5em}

\pagestyle{fancy}
\fancyhf{}
\lhead{\small VV-F -- Vecný zámer projektu}
\rhead{\small Akronym: BlackBox}
\cfoot{\small \thepage{} / \pageref{LastPage}}

\begin{document}

\begin{center}
  {\Large\bfseries VV-F -- VECNÝ ZÁMER PROJEKTU}\\[0.3em]
  Záväzná osnova pre aplikovaný výskum a~vývoj (maximálny rozsah 15 strán)
\end{center}

\renewcommand{\arraystretch}{1.3}
\noindent
\begin{tabularx}{\linewidth}{|>{\bfseries}p{45mm}|X|}
\hline
Názov projektu & Black Box: lokálna aplikácia na včasné varovanie študentov pred študijným preťažením \\ \hline
Zodpovedný riešiteľ & Tymofii Trembovetskyi \\ \hline
Žiadateľ & Slovenská technická univerzita v~Bratislave, Fakulta informatiky a~informačných technológií \\ \hline
Štatutárny zástupca žiadateľa & prof. Dr. Ing. Oliver Moravčík, rektor \\ \hline
\end{tabularx}
\renewcommand{\arraystretch}{1}

%=====================================================================
\section{Excelentnosť}

\subsection{Projektový zámer, spôsob riešenia a~aplikačná úroveň výsledkov}

Študenti dostávajú termíny zadaní, testov a~skúšok z~informačného systému
školy, e-mailov a~skupinových chatov. Každý zdroj osve vyzerá zvládnuteľne,
no keď sa termíny z~viacerých predmetov zbehnú do jedného týždňa, študent to
zistí až niekoľko dní vopred. Opakované obdobia preťaženia sú pritom jednou
z~ciest k~vyhoreniu~\cite{Maslach:Burnout,Schaufeli:Burnout}.

Projektovým zámerom je navrhnúť a~overiť desktopovú aplikáciu Black Box,
ktorá:
\begin{enumerate}[nosep]
  \item zhromaždí termíny na jednom mieste (ručné zadanie, import vo formáte
    iCalendar~\cite{RFC5545}),
  \item z~odhadov časovej náročnosti termínov vypočíta záťaž jednotlivých
    týždňov a~7 až 14 dní vopred varuje pred týždňom, ktorý bude
    pravdepodobne preťažený,
  \item umožní jednoduchý denný záznam pocitov a~po náročnom období zobrazí
    spoločnú časovú os termínov, záťaže a~pocitov.
\end{enumerate}

Riešenie stojí na troch pilieroch. Prvým je \emph{model záťaže}, ktorý
rozkladá odhad hodín do dní pred termínom s~rastúcou váhou a~porovnáva
týždennú záťaž s~kapacitou študenta aj s~jeho zvyčajnou záťažou. Druhým je
\emph{kalibrácia odhadov}, ktorá koriguje systematické podceňovanie času.
Tretím je architektúra \emph{local-first}~\cite{Kleppmann:LocalFirst}:
aplikácia nemá serverovú časť a~citlivé údaje o~pohode nikdy neopustia
zariadenie používateľa.

Aplikačnou úrovňou výsledkov je funkčný prototyp overený so študentmi
počas reálneho semestra (TRL~6), ktorý je možné bez ďalšieho vývoja
ponúknuť študentom fakulty a~ďalších vysokých škôl.

\subsection{Aktuálnosť cieľov a~originalita}

Duševné zdravie študentov je aktuálna téma: medzinárodný prieskum
Svetovej zdravotníckej organizácie zistil, že približne tretina študentov
prvého ročníka mala v~poslednom roku príznaky niektorej z~bežných duševných
porúch~\cite{Auerbach:WMH}, a~akademický stres patrí medzi ich hlavné
zdroje~\cite{Pascoe:Stress}. Zároveň rastie počet aplikácií na sledovanie
pohody, z~ktorých veľká časť odosiela údaje tretím
stranám~\cite{Huckvale:Privacy}.

Originalita myšlienky spočíva v~prepojení troch oblastí, ktoré sa doteraz
riešili oddelene: správy termínov, sledovania pohody a~včasného varovania.
Originalita riešenia spočíva v~prenesení princípu pomeru akútnej
a~chronickej záťaže zo športovej vedy~\cite{Gabbett:ACWR} do študijného
prostredia, jeho úprave podľa známej kritiky~\cite{Impellizzeri:ACWR}
a~doplnení o~kalibráciu odhadov, ktorá reaguje na plánovací
klam~\cite{Buehler:Planning}. Nepoznáme nástroj, ktorý by tieto prvky
spájal a~zároveň uchovával všetky údaje lokálne.

\subsection{Stav poznania, vedecké základy a~predchádzajúce výsledky}

Podrobný prehľad stavu poznania prináša sprievodný prehľadový článok;
tu uvádzame jeho hlavné body.

\paragraph{Záťaž a~jej vnímanie.}
Vyhorenie študentov je dobre preskúmaný a~merateľný
jav~\cite{Schaufeli:Burnout}. Vnímaná záťaž nezodpovedá jednoducho počtu
hodín učenia, ale závisí aj od organizácie výučby a~hodnotenia~\cite{Kember:Workload}.
Študenti systematicky podceňujú čas potrebný na úlohy~\cite{Buehler:Planning}
a~sústreďujú úsilie do obdobia tesne pred termínom~\cite{Steel:Procrastination},
preto sa prekryv termínov prejaví silnejšie, než by naznačoval rovnomerný plán.

\paragraph{Sledovanie pohody.}
Projekt StudentLife ukázal, že stres študentov rastie a~spánok aj aktivita
klesajú v~obdobiach so zvýšeným počtom termínov~\cite{Wang:StudentLife}.
Metóda ekologického momentového hodnotenia (EMA) poskytuje spoľahlivejší
obraz pohody ako spätné dotazníky~\cite{Shiffman:EMA}. Model osobnej
informatiky upozorňuje, že používatelia zlyhávajú najmä pri spájaní údajov
z~rôznych zdrojov a~pri reflexii~\cite{Li:PersonalInformatics}.

\paragraph{Včasné varovanie.}
Systémy učiacej sa analytiky, napríklad Course
Signals~\cite{Arnold:CourseSignals}, varujú pred neúspechom v~predmete, sú
však inštitucionálne a~nesledujú záťaž naprieč predmetmi. V~športovej vede sa
riziko preťaženia hodnotí pomerom akútnej a~chronickej
záťaže~\cite{Gabbett:ACWR}; jeho prahy však nie sú prenosné bez overenia
v~konkrétnom kontexte~\cite{Impellizzeri:ACWR}.

\paragraph{Súkromie.}
Údaje o~pohode môžu byť údajmi o~zdraví, ktoré GDPR radí medzi osobitné
kategórie~\cite{GDPR,Zakon18}. Prístup local-first umožňuje poskytnúť plnú
funkčnosť bez toho, aby údaje opustili zariadenie~\cite{Kleppmann:LocalFirst}.

\paragraph{Predchádzajúce výsledky riešiteľov.}
Riešitelia vypracovali projektový zámer, prehľadový článok a~prvý návrh
modelu záťaže, ktorý je opísaný v~časti~\ref{s:metodologia}.

\subsection{Ciele projektu a~reálnosť ich dosiahnutia}

\textbf{Hlavný cieľ:} navrhnúť, implementovať a~overiť lokálnu aplikáciu,
ktorá študentovi vopred ukáže preťažené týždne semestra a~po ich uplynutí
pomôže pochopiť ich príčiny.

\begin{center}
\small
\begin{tabularx}{\linewidth}{@{}lXX@{}}
\toprule
& \textbf{Čiastkový cieľ} & \textbf{Merateľný ukazovateľ} \\
\midrule
C1 & Model týždennej záťaže a~pravidlo varovania
   & presnosť varovaní $\geq 70\,\%$, úplnosť $\geq 60\,\%$ pri predstihu
     $\geq 7$ dní; vyššie F1 ako východiskový model \\
C2 & Kalibrácia odhadov časovej náročnosti
   & pokles strednej absolútnej chyby odhadu po kalibrácii \\
C3 & Prototyp aplikácie pre Windows, Linux a~macOS
   & $\geq 95\,\%$ správne načítaných udalostí iCalendar; žiadna
     odchádzajúca komunikácia s~osobnými údajmi \\
C4 & Overenie so študentmi počas semestra
   & $\geq 30$ účastníkov, 8 týždňov; skóre SUS $\geq 71$~\cite{Bangor:SUS};
     denný záznam v~$\geq 50\,\%$ dní \\
\bottomrule
\end{tabularx}
\end{center}

\textbf{Reálnosť dosiahnutia.} Ciele C1 a~C2 sa opierajú o~jednoduché
a~dobre interpretovateľné výpočty, ktoré je možné overiť najprv na
simulovaných rozvrhoch. Cieľ C3 využíva zrelé technológie (Tauri~\cite{Tauri},
SQLite~\cite{SQLite}) a~štandardizovaný formát iCalendar, ktorý exportuje
aj akademický informačný systém STU. Pre cieľ C4 je k~dispozícii veľká
skupina potenciálnych účastníkov na vlastnej fakulte. Hodnoty ukazovateľov
v~C1 sú zámerne stanovené ako dosiahnuteľné minimum; ak ich model nedosiahne,
bude to samo osebe dôležitým zistením o~predvídateľnosti preťaženia.

\subsection{Metodológia riešenia} \label{s:metodologia}

\paragraph{Model záťaže (C1).}
Pre termín $i$ s~dátumom $d_i$, odhadom hodín $h_i$ a~dĺžkou prípravy
$w_i$ dní je záťaž dňa $t$
\begin{equation*}
L(t) = \alpha \sum_i h_i \, k_i(d_i - t), \qquad
k_i(x) = \frac{2\,(w_i - x)}{w_i\,(w_i + 1)}, \quad x = 0, \ldots, w_i - 1,
\end{equation*}
a~záťaž týždňa $T$ je $Z(T) = \sum_{t \in T} L(t)$. Lineárne jadro
$k_i$ dáva väčšiu váhu dňom tesne pred termínom, čo zodpovedá poznatkom
o~prokrastinácii~\cite{Steel:Procrastination}. Varovanie sa vydá, ak
$Z(T) > c$ (kapacita študenta) alebo ak pomer $Z(T)$ k~priemeru štyroch
predchádzajúcich týždňov (s~dolnou hranicou $c/2$) prekročí $\rho$.
Počiatočné hodnoty $c = 20$~h a~$\rho = 1{,}5$ budú v~etape E1 nastavené
na simulovaných rozvrhoch a~v~etape E3 overené na reálnych údajoch.
Výber jednoduchého interpretovateľného modelu namiesto strojového učenia
je zámerný: aplikácia musí vedieť študentovi vysvetliť, prečo varuje,
a~pri lokálnom spracovaní nemá k~dispozícii údaje iných používateľov.

\paragraph{Kalibrácia (C2).}
Po termíne aplikácia voliteľne zistí skutočný počet hodín $s_i$.
Koeficient $\alpha$ je medián pomerov $s_i / h_i$ za posledných najviac
desať termínov (pri menej ako piatich termínoch $\alpha = 1$).

\paragraph{Vývoj (C3).}
Aplikácia sa vyvíja iteratívne v~dvojtýždňových iteráciách. Jadro
v~jazyku Rust obsahuje model, kalibráciu a~prístup k~databáze a~je pokryté
jednotkovými testami; import iCalendar sa testuje na anonymizovaných
exportoch z~informačného systému. Používateľské rozhranie sa overuje
priebežne s~piatimi až siedmimi študentmi metódou myslenia nahlas.

\paragraph{Overenie (C4).}
Pilotné overenie prebehne počas letného semestra 2028 s~najmenej 30
dobrovoľníkmi z~rôznych ročníkov a~študijných programov. Účastníci
podpíšu informovaný súhlas. Raz týždenne odpovedajú na otázku, či bol
uplynulý týždeň preťažený; týždeň sa za preťažený považuje aj vtedy, ak je
priemer denného záznamu najviac 2. Na konci overenia účastník sám
exportuje anonymizovaný súhrn ukazovateľov (bez názvov predmetov
a~termínov), ktorý odovzdá riešiteľom. Vyhodnotí sa presnosť, úplnosť
a~F1 varovaní oproti východiskovému modelu (počet termínov v~týždni),
pokles chyby odhadu po kalibrácii (párový Wilcoxonov test) a~skóre
SUS~\cite{Brooke:SUS}.

\subsection{Najdôležitejšie výstupy zodpovedného riešiteľa za posledných 5~rokov}

Zodpovedný riešiteľ je študentom bakalárskeho štúdia na FIIT STU
a~doteraz nemá publikované výstupy aplikovaného výskumu. Prvým výstupom
v~tejto oblasti je prehľadový článok, ktorý je súčasťou tohto návrhu.

\subsection{Najdôležitejšie projekty zodpovedného riešiteľa za posledných 5~rokov}

Zodpovedný riešiteľ doteraz neriešil projekt financovaný z~grantovej schémy.

\subsection{Úroveň kvality zodpovedného riešiteľa, jeho vízia a~kariérny rast}

Zodpovedný riešiteľ študuje informatiku a~tému pozná aj z~vlastnej
skúsenosti študenta. Projekt pre neho predstavuje prvú samostatnú
výskumnú prácu, v~ktorej si osvojí celý cyklus aplikovaného výskumu: od
prehľadu literatúry cez návrh a~implementáciu po overenie s~používateľmi.
Víziou je pokračovať v~téme v~bakalárskej a~diplomovej práci, rozšíriť
model o~prispôsobovanie prahov jednotlivcovi a~zapojiť do vývoja ďalších
študentov ako prispievateľov do otvoreného zdrojového kódu.

\subsection{Kompetentnosť riešiteľskej organizácie}

Fakulta informatiky a~informačných technológií STU sa dlhodobo venuje
softvérovému inžinierstvu, interakcii človeka s~počítačom a~analýze
údajov o~učení. Fakulta má skúsenosti s~vedením študentských výskumných
projektov a~poskytuje prostredie, v~ktorom je možné pilotné overenie
uskutočniť priamo medzi jej študentmi.

\subsection{Expertíza riešiteľov}

Kyryl Tkachenko bude zodpovedný za používateľské rozhranie, návrh
denného záznamu a~časovej osi. Alex Vyšný bude zodpovedný za organizáciu
pilotného overenia, nábor účastníkov a~štatistické vyhodnotenie údajov.
Zodpovedný riešiteľ bude zodpovedný za jadro aplikácie, model záťaže
a~kalibráciu. Všetci riešitelia sa podieľajú na vyhodnotení a~publikovaní
výsledkov.

\subsection{Zapojenie mladých pracovníkov}

Všetci traja riešitelia sú mladí pracovníci do 35 rokov, zapojenie mladých
pracovníkov je teda 100\,\% kapacity projektu. Projekt je pre nich
príležitosťou získať skúsenosti s~výskumom, ktoré môžu zúročiť
v~ďalšom štúdiu.

%=====================================================================
\section{Dopad}

\subsection{Prínosy pre rozvoj poznania, aplikovaného výskumu a~inovácií}

\begin{itemize}[nosep]
  \item \textbf{Nová služba:} aplikácia, ktorá študentovi vopred ukáže
    preťažené týždne, čo dnes žiadny bežne dostupný nástroj neumožňuje.
  \item \textbf{Nový postup:} model záťaže s~kalibráciou, prenesený zo
    športovej vedy a~upravený pre študijné prostredie, spolu s~metodikou
    overenia jeho presnosti.
  \item \textbf{Poznatky:} empirické údaje o~tom, nakoľko je študijné
    preťaženie predvídateľné z~termínov a~odhadov a~ako veľmi študenti
    podceňujú časovú náročnosť úloh.
  \item \textbf{Spoločenská inovácia:} ukážka, že nástroj pracujúci
    s~citlivými údajmi o~pohode môže byť užitočný bez toho, aby údaje
    opustili zariadenie.
\end{itemize}

\subsection{Využiteľnosť výsledkov v~SR a~v~zahraničí}

Aplikácia bude zverejnená pod otvorenou licenciou a~bude ju môcť bezplatne
používať ktorýkoľvek študent. Žiadateľ ju môže ponúknuť študentom prvého
ročníka v~rámci adaptačných aktivít; zistenia o~rozložení záťaže môžu
garanti predmetov využiť pri plánovaní termínov. Keďže formát iCalendar
je medzinárodný štandard, aplikácia je použiteľná aj na iných
univerzitách v~SR a~v~zahraničí. Rozhranie bude dostupné v~slovenčine
a~angličtine.

\subsection{Ekonomický a~spoločenský prínos}

Hlavný prínos je v~oblasti kvality života a~zdravia: včasné varovanie
umožňuje rozložiť prácu skôr, než vznikne stres, a~znižuje tak riziko
dlhodobého preťaženia a~vyhorenia~\cite{Maslach:Burnout}. Akademický
stres súvisí s~horším študijným výkonom a~predčasným ukončovaním
štúdia~\cite{Pascoe:Stress}; každý študent, ktorý štúdium vďaka lepšiemu
plánovaniu dokončí, predstavuje úsporu verejných prostriedkov vynaložených
na jeho vzdelávanie. Projekt zároveň skvalitňuje ľudské zdroje: riešitelia
získajú skúsenosti s~aplikovaným výskumom.

\subsection{Opatrenia na maximalizáciu výsledkov a~komunikácia}

\begin{itemize}[nosep]
  \item zverejnenie zdrojového kódu a~inštalačných balíkov vo verejnom
    repozitári s~dokumentáciou,
  \item webová stránka projektu s~návodom a~priebežnými výsledkami,
  \item príspevok na vedeckej konferencii a~prehľadový článok,
  \item workshop pre študentov o~plánovaní záťaže s~ukážkou aplikácie,
  \item prezentácia výsledkov vedeniu fakulty a~študijnej rade.
\end{itemize}

%=====================================================================
\section{Implementácia}

\subsection{Projektový plán a~výstupy}

Projekt trvá 24 mesiacov (01/2027 -- 12/2028) a~je rozdelený do štyroch
pracovných balíkov (tab.~\ref{t:balicky}, obr.~\ref{f:gantt}).

\begin{table}[htb]
\centering
\small
\caption{Pracovné balíky, výstupy a~míľniky.}
\label{t:balicky}
\begin{tabularx}{\linewidth}{@{}l>{\raggedright\arraybackslash}p{40mm}X@{}}
\toprule
& \textbf{Pracovný balík} & \textbf{Hlavné úlohy a~výstupy} \\
\midrule
E1 & Analýza a~model záťaže (M1--M6)
   & analýza exportov z~informačného systému, návrh modelu a~kalibrácie,
     simulátor; \emph{výstup:} technická správa, simulátor;
     \emph{míľnik M1 (M6):} model overený na simulovaných rozvrhoch \\
E2 & Vývoj prototypu (M4--M13)
   & jadro v~Ruste, databáza, import iCalendar, rozhranie, varovania,
     denný záznam, časová os; \emph{výstupy:} verzia 0.1 (M9), verzia 1.0
     (M13); \emph{míľniky M2 (M9), M3 (M13)} \\
E3 & Pilotné overenie (M14--M19)
   & nábor účastníkov, informovaný súhlas, 8-týždňové používanie,
     zber súhrnov; \emph{výstup:} správa o~overení;
     \emph{míľnik M4 (M19):} ukončený zber údajov \\
E4 & Vyhodnotenie a~šírenie (M18--M24)
   & analýza údajov, úprava modelu, verzia 2.0, publikácia, workshop;
     \emph{míľnik M5 (M24):} záverečná správa \\
\bottomrule
\end{tabularx}
\end{table}

\begin{figure}[htb]
\centering
\begin{tikzpicture}[x=6mm, y=6mm, font=\small]
  \foreach \m in {1,...,24} \node[font=\scriptsize] at (\m-0.5, 0.6) {\m};
  \draw[gray!40] (0,0) -- (24,0);
  \foreach \row/\a/\b/\n in {1/0/6/E1, 2/3/13/E2, 3/13/19/E3, 4/17/24/E4} {
    \node[anchor=east] at (-0.2,-\row+0.5) {\n};
    \fill[black!25] (\a,-\row+0.15) rectangle (\b,-\row+0.85);
  }
  \foreach \x/\l in {6/M1, 9/M2, 13/M3, 19/M4, 24/M5} {
    \draw[dashed, gray] (\x,0) -- (\x,-4.2);
    \node[font=\scriptsize, below] at (\x,-4.2) {\l};
  }
\end{tikzpicture}
\caption{Harmonogram pracovných balíkov (mesiace projektu) a~míľniky.}
\label{f:gantt}
\end{figure}

\subsection{Projektový manažment}

Projekt riadi zodpovedný riešiteľ. Riešitelia sa stretávajú raz týždenne;
na konci každej dvojtýždňovej iterácie vyhodnotia splnené úlohy
a~naplánujú ďalšie. Úlohy, zdrojový kód a~dokumentácia sú vedené
v~systéme na správu verzií s~evidenciou úloh. Na konci každej etapy
riešitelia porovnajú stav s~míľnikmi a~ukazovateľmi z~časti~1.4 a~v~prípade
odchýlky upravia plán ďalšej etapy.

\subsection{Riziká implementácie a~ich zmiernenie}

\begin{center}
\small
\begin{tabularx}{\linewidth}{@{}>{\raggedright\arraybackslash}p{45mm}cX@{}}
\toprule
\textbf{Riziko} & \textbf{Pravdep.} & \textbf{Zmiernenie} \\
\midrule
Nedostatok účastníkov overenia & stredná
  & nábor vo viacerých ročníkoch, odmena vo forme poukážky, rezerva 10
    účastníkov nad minimum \\
Účastníci prestanú zapisovať údaje & vysoká
  & záznam jedným kliknutím, pripomienka, sledovanie podielu vyplnených dní
    už počas overenia \\
Model nedosiahne cieľovú presnosť & stredná
  & nastavenie prahov na simulovaných údajoch v~E1, prispôsobenie prahov
    jednotlivcovi v~E4; negatívny výsledok sa zverejní \\
Export informačného systému nie je kompatibilný & nízka
  & testy na reálnych exportoch v~E1, ručné zadanie ako záložná možnosť \\
Únik citlivých údajov & nízka
  & žiadna serverová časť, overenie sieťovej komunikácie, anonymizovaný
    súhrn exportuje sám účastník \\
\bottomrule
\end{tabularx}
\end{center}

\subsection{Adekvátnosť rozpočtu}

Celkový rozpočet je 63\,274~EUR (formulár VV-C). Najväčšiu časť tvoria
osobné náklady troch riešiteľov (2\,400 hodín za projekt), ktoré
zodpovedajú rozsahu vývoja a~overenia. Materiál zahŕňa dva testovacie
notebooky, bez ktorých nie je možné zostaviť a~otestovať aplikáciu pre
macOS a~Windows. Služby pokrývajú certifikáty na podpisovanie kódu,
odmeny účastníkom overenia a~konferenčný poplatok. Projekt nepožaduje
nákup prístrojov ani licencií, keďže všetky použité technológie sú
otvorené.

\subsection{Existujúca infraštruktúra pracoviska}

FIIT STU poskytuje pracovné priestory, počítačové učebne, prístup
k~sieti a~výpočtovým prostriedkom a~licencie na bežný softvér.
Riešitelia majú k~dispozícii vlastné počítače s~operačným systémom
Linux, na ktorých prebieha vývoj. Pilotné overenie nevyžaduje žiadnu
ďalšiu infraštruktúru, pretože aplikácia beží výlučne na zariadeniach
účastníkov.

%=====================================================================
\bibliographystyle{plain}
\bibliography{literatura}

\end{document}
